\clearpage
\SourceLeaf{39}

\begin{center}
\textbf{[MECHANISCHER SATZ ÜBER DIE WURZELN EINER GANZEN
FUNCTION $f(x)$ UND IHRER ABLEITUNG $f'(x)$.]}
\end{center}

\noindent\hspace*{\parindent}Sind $\begin{array}{c}a,\,b,\,c,\,.\,.,\,m,\,n\\ a',\,b',\,c',\,.\,.,\,m'\end{array}$ die
Wurzeln der Gleichung $\begin{array}{c}f\,x=0\\ f'\,x=0,\end{array}$ wo $f'x=\dfrac{d\,f\,x}{d\,x}$,
und werden durch dieselben Buchstaben die entsprechenden Punkte in plano bezeichnet, so ist, wenn
man sich in $a,\,b,\,c,\,.\,.,\,m,\,n$ gleiche abstossende oder anziehende Massen denkt, die im
umgekehrten Verhältniss der Entfernung wirken, in $a',\,b',\,c',\,.\,.,\,m'$ Gleichgewicht.

\Needspace{9\baselineskip}
\begin{center}
\textbf{BEMERKUNGEN ZU DEN BEIDEN VORSTEHENDEN ALGEBRAISCHEN NOTIZEN.}
\end{center}

\noindent\hspace*{\parindent}Die erste der beiden mitgetheilten Notizen stammt wahrscheinlich aus den Jahren 1826--28. Sie
bezweckt die Auflösung eines Systems von fünf linearen Gleichungen mit besonders einfach und zwar
symmetrisch gebauter Determinante vermittelst der fünften Einheitswurzeln. Beim Beweise der
schliesslich herauskommenden Darstellungen für die fünf Grössen $p,\,q,\,r,\,s,\,t$ beachte man,
dass die Norm der complexen Zahl $a+b\varepsilon+c\varepsilon\varepsilon+d\varepsilon^3+e\varepsilon^4$
sich in der Gestalt $\frac{1}{4}(5M-m\mathfrak{M})$ darstellen lässt. 

Die zweite Notiz hat \textsc{Gauss} wahrscheinlich im Jahre 1846 aufgezeichnet. Zum Beweise
orientire man die Ebene so, dass eine der Lösungen von $f'(x)=0$ in den Nullpunkt zu liegen kommt.
Dann ist:

\[
\begin{aligned}
\frac{1}{a}+\frac{1}{b}+\frac{1}{c}+\cdots+\frac{1}{n}&=0,\\[3pt]
\frac{a_0}{a\,a_0}+\frac{b_0}{b\,b_0}+\frac{c_0}{c\,c_0}+\cdots+\frac{n_0}{n\,n_0}&=0,
\end{aligned}
\]

\noindent wo $a_0$ zu $a$, $b_0$ zu $b$, etc. conjugirt complex sein sollen. Die mechanische Deutung der
letzten Gleichung liefert den \textsc{Gauss}’schen Satz.
\par\nobreak\hfill\textsc{Fricke.}\par

